Understanding when a system is correct, secure, and stable has
traditionally been addressed by three largely independent lines of work:
logical and formal systems study correctness via fixed points;
security and adversarial models analyze robustness via attack surfaces;
game-theoretic and dynamical systems investigate stability via equilibria.

\subsection{Main result: a conditional collapse}

We show that correctness, unbreakability, and stability collapse under:
(i) certificate coverage (C6); (ii) no unavoidable external inputs; (iii) unique attractor.

This equivalence is not assumed---it arises from the elimination of all degrees of freedom.

\subsection{Degrees of freedom as a unifying invariant}

Three types: representational $\DoF$ (alternative admissible states);
source $\DoF = 1/\Cmin$ (finite-cost break paths);
dynamic $\DoF$ (alternative stable behaviors).

$\DoF(P)=(0,0,0)$ is the exact condition for three-way equivalence.

\subsection{Contributions}

\begin{enumerate}[label=(\arabic*)]
  \item Minimal unified language $(\At,F,\mathcal{B},\mathcal{A})$.
  \item DoF Classification Completeness: no fourth DoF exists (Theorem~\ref{thm:dof-completeness}).
  \item Synthesis Theorem: $\DoF=(0,0,0) \iff$ three-way equivalence.
  \item Boundary: $\UES(P)\neq\emptyset$ is the precise cause of failure in Class III.
  \item Execution Layer: $C_{\min}^{ctrl}$ attained at execution layer (Theorem~\ref{thm:el}).
  \item Empirical: all three DeFi events are $\DoF_{\mathrm{src}}>0$ without adversary.
\end{enumerate}

\subsection{Organization}

\S\ref{sec:core}: minimal core. \S\ref{sec:expressivity}: expressivity.
\S\ref{sec:synthesis-main}: DoF and Synthesis Theorem.
\S\ref{sec:cot}: boundary. \S\ref{sec:mlr}: open systems.
\S\ref{sec:defi-empirical}: DeFi verification.
